\documentclass{article}
\usepackage[margin=0.8in]{geometry}
\usepackage{hyperref}

\title{Cornelius Lanczos}
\author{The University of Manchester}
\date{1972}

\begin{document}
\maketitle

\newpage
\tableofcontents
\newpage

\section{About His Life}
\subsection{Introduction}
Ladies and gentlemen, I consider it a real privilege indeed that the videotape is made of an entirely personal nature and that I have a chance to talk to you today in the form of an informal chat about my own life. I don't do that at all in order to show you what a wonderful man I am on the other hand, I want to stay away from any false modesty either. I consider the value of such a enterprise in the fact that I had the privilege of having a long life and a varied life. I encountered, in my life, a great number of eminently interesting people, excellent people in various countries and in various situations. I had a chance to be exposed not only to excellent people and excellent situations but to great ideas and contribute myself to some extent to the development of great ideas. Therefore, there is quite a bit what I can say. naturally I would like to say things always in such a way that I put it in a broader background because it is not my own personality which is important but it is the interrelation between the ego \textit{I} and the surrounding world which I think might be of interest.

\subsection{Childhood and education in Hungary}
I start out with the fact that I was born before the beginning of this century. I was born in 1893 in Hungary in a village with about 30,000 inhabitants and about 60 kilometers or 40 miles to the southwest from Budapest. It was a town of remarkable vigour and remarkable intellectual interest where people lived their life which was rich in many respects; not in material goods but rich in intellectual values. One probably has a wrong impression on what such a small country as hungary actually was. But in a sense it was a gem in the line of gracious living. Gracious living in the sense that one consider it a privilege to study and to participate in intellectual endeavors; in music, art, and science. So it was very different from the kind of permissive society we have today. When I grew up it was the late Victorian era by which one could characterize the middle European scene of those days. The parents were on a very high pedestal and were highly removed from a person's actual presence. One didn't discuss any kind of intimate matters with the parents and one had certain duties automatically in connection with privileges. 

It was not the idea that certain education that is coming to me we considered it a privilege because after all the educated middle class that was already a the kind of capitalistic selection granted. It was it was a selection. It was not everybody's job to be able to study and therefore we appreciated the fact that we could study and that it was automatically our duty to express our appreciation by working hard. The so called secondary school education was on a high level. It was not the question of what one learned but one developed certain attitudes which were tremendously important because they accompany the person throughout his life. The study was a duty and also a privilege. We found it wonderful to be exposed to all those ideas through which the human race went and since at that time in these pre-puberty years one was still free of that emotional upset which came later on these years spent in the secondary school were formative years of very great importance to a person's life. As I told you we had of course all subjects in sciences arts languages wasn't so called humanistic gymnasium where Latin and Greek was thought I want one if one wanted one would be or could be free of their education in Greek but Latin was compulsory and I think we had it for eight years but it's a wonderful wonderful background a classical background I think is very important and very good and so these years very important the opened doors the open channels in so many directions and we took advantage of that we had discussions and we had a club where the students themselves came together and had discussions in all in all all kind of adventures some of them were interested in poetry some of some of them in art some of them in sciences and by presenting essays and presenting talks to the rest of the students we came already into a frame of mind where we saw how variable discussions are and that however should stay away from a kind of dogmatic attitude by the fact that the same issue has various aspects and that one can discuss under various as aspects I consider that a very important experience to stay away from dogmatism and one-sidedness by seeing that the same issue has various aspects we were used to that from our secondary days on secondary school year days on now when I finished in 1910 made my so-called maturity examination which was the entrance business in a sense the entrance examination to the University of Man was automatically taken and at the university I started do I developed from my about 15 year on I had a very strong pre-election toward the sciences to add mathematics physics particularly mathematics and so it was for me no question that I wanted to pursue mathematics as a profession very much against the wish of my father who was an eminent lawyer a very excellent man highly appreciated in that town but only enough somewhat pragmatic mind who said you have to look out you have to look out in life what are you going to do now what can you do with mathematics there were no industrial labs in Hungary there were no one couldn't count on a university career because they were so exceedingly few positions free and so what one could do is you prepared to become a secondary school teacher of mathematics or physics and that happened to be an occupation which didn't appeal to my father at all however I have been able to break down his antagonism and

\subsection{Student in Budapest}
started at the University of Budapest in 1911 to study mathematics physics and as a side subject philosophy this philosophy course meant a tremendous tremendously much to me because it was an outstanding expository I mean this man Bernard Alexander who was an outstanding man as an expository a teacher he was a born teacher he was not a man who did very much in the history of philosophy are as at was his subject history philosophy he wrote a few articles in the book on Spinoza and a few and the or country was a tremendous expert of content but what was his strength was that he put himself completely I mean starting with the old Greeks and coming up to modern times he always put himself completely in the identified himself completely with the philosopher whom he just happened to discuss so that when he talked on to us it looked that everything is just wonderful it must be that this is the philosophy then when he came then to fished it it turned out everything was wrong that has done and so his his idea was what he called eminent eminent criticism one should discuss one should criticize a system first of all from his own standpoint and then one can criticize it from from the standpoint of an outsider then monkey creatures are criticized from the standpoint of how does it fit into the general frame of philosophy but first of all you have to understand it from its own standpoint how consistent it is what can you do with it how much does it help you to understand the great problems of philosophy this was to me a wonderful experience and I enjoyed his course so much that I think I am I took it twice this course in the history of philosophy because it was such a great experience in addition to that I took the usual courses in mathematics and physics and it happened that one didn't go only to the universe typical there was also in of course I spent these years in Budapest I got away from home and I spent these four or five years until ending because it was not only the four years University course you had to finish also you have to take the so-called a pedagogical or an educational courses in order to be privileged in order to get the license for teaching so it was approximately five years before you really finished and because you had to have also the courses in the history way in educational matters and you had to have also some practice teaching and so on now I was of course lucky because I did make a university career from the beginning and my father was more than happy and more than satisfied when he saw that I didn't have to go to a secondary school and for the rest of my life and

\subsection{Emigration to Germany, research assistant in Friburg and Frankfurt}
I graduated then approximately in 1902 15 16 I started to work on my research problems I wanted to make PhD but I've got or an appointment as an assistant at the Polytechnic or as I told you the university was not the only place one also went to the technical engineering school is Politecnico which was also an institution of very high standing and took certain courses in mathematics physics there so that I was acquainted with the staff of both of these institutions and I was when I when I graduated after my graduation I got an appointment to become an assistant of physics in the Politecnico I spent there four years from 916 to 20 these were bad years in a certain sense because Hungary was unfortunately the first country in Europe where the ugly head of anti-semitism raised itself already during the war and particularly after the war the propaganda was that the Jews did not do enough that not enough Jews died in the world that even privileged positions and therefore they tried to dodge the military service personally I was not taken to them to military service on account of my physical makeup which was considered unsuited for military service but the actual statistics shows that this is absolutely untrue the Jews died like anybody else at least in the proportion of of the popular population so that there was absolutely no truth in the alligator alligator that they were not patriotic they were not sufficiently patriotic and that they wanted to dodge the voyage but it started already during the war and after the war there was a very unfortunate episode for half a year Hungary was under a communistic is being béla kun a communistic Soviet a type of government in which unfortunately there were certain Jews in the government now the revenge was terrible it came a kind of in the form of a white terror so-called white terror which was just as bad as a NASA terror only it didn't last very long because the Hungarian people are not very subject to prolonged hatreds and although you can you can certainly buy propaganda you can get them to do anything but it doesn't hold out for very long so this extreme fascistic honestok type of hate propaganda in which thousands of jews died and one had to be one had to watch out for a perform for one's life did not last very long for about half a year however when my I mean I personally had absolutely no difficulty zombies never I got along very well with the staff I got along with the students very well I had no personal difficulties everybody knew that I'm a Jew and I was never hiding that and there was no difficulty and when these four years were over then my chair and my boss told me that I'm sorry but I am not in the position to renew your you know what the conditions are I cannot renew your contract so I resigned in 920 and prepared for my PhD I got my PhD at the University of sigit which was a country university in budapest away from budapest i got it there from a professor with whom i was on very good terms he was not a very productive mind but exceedingly exceedingly good-hearted and exceedingly taking taking the issue of the young physicists study the great deal for Hungarian physics and I made he made me a PhD degree with him in physics as a major subject mathematics is a second major subject philosophy as a minor subject so that brings me now to the year of 1921 on account of the so-called Jewish law there was no which restricted the in of course we must not forget the Jews did play a very prominent role in the professionals in the professions in the lawyers the doctors the journalists we are more than 50 in Budapest of the Jewish extraction so that there was there was a certain that this must have caused a certain jealousy that was obvious then came the Jewish law which restricted all these activities to the percentage of the population which was only 7 which meant that the Jews were practically cut off of the higher studies because they were not taken at the university anymore

\subsection{Assistant of Albert Einstein in Berlin, antisemitism in Germany}
this had in a certain sense very beneficial effects because there was a there was a phase difference the entire Semitism in Hungary started around 1920 around these years in Germany the same movement started 33 but there were these 13 years between so some outstanding hungarians is for example vigener and Ronnie for Norman teller Silla quite a number of outstanding physicists who left the country because we are forced to leave the country went to Germany and in in those days you could study in Germany and you were fully accepted and these exceptionally gifted people came in two positions which correspond to their gifts and they are the chance to do the Graduate Studies there and come into the proper positions when Hitler came these people were already so famous in their variety so well established in their sciences that the United States recognized the importance of getting such outstanding people and they immediately got a position in the United States so it worked out this phase difference between the anti-semitic excesses in Hungary and later in Germany worked out exceedingly well for the intellectual elite that there were many victims undoubtedly but as far as the intellectual elite goes it worked out very well and these people came to Germany then and from Germany they came to the United States to some extent to England but to a large extent to the United States and I personally in 21 made my after making my PhD degree I made an application to a German university and I was taken immediately because to my surprise the fact that I was a Jew didn't play any role but after all I was a foreigner from Hungary and some of these new German universities were pretty much naturalistically minded but it was the deepest inflation in Germany a Germany was in the clutches of the inflation when when the mark went into the millions and billions and lost all value at that time anybody who amounted to something any physicist or engineer are basically also the scientifically trained people worked for industry because there they are they're paid in dollars as ever they were paid in hard currency and so they they were not lost and in the university positions were free he advertised to get people I got immediately a position in Freiburg I'm Breisgau I enjoyed it tremendously I made very good contacts with movies my boss who was a German nationalist of the old school very conservative in his outlook but a charming man and we got along very well he became seventy years at that time and a new successor brought his assistants along so I had to look out for something else I got immediately another position in Frankfurt in mind that was in 1924 I went to Frankfurt I stayed there for seven years up to 31 and I published already at that time I started to publish in the field of relativity which attracted me tremendously it was something of a tremendous experience to come in touch with the Einsteinian ideas which at first as a young undergraduate student I rejected as absurd as so many other people do but when I really understood what he wants to do and what these ideas are are that was like like a champagne I mean that was something tremendously exciting and I started to work I started to publish already in 94 vacation came out in 1922 and I had a number of publications in connection with Einstein and gravitational theory and it so happened that Einstein in the year of 1919 27 he had an assistant who Dahmer who for some reasons he had to leave and he was looking looking for somebody and Sylla who worked with him at that time he was influential he knew my papers to suggest me as his assistant the city thus it happened that in 92 1928 29 I got a leave of absence from the university of frankfurt to work in berlin now this year was a tremendous year because to see such a great man almost daily and have a chance to talk to him about mutually interesting problems in physics was of course an enormous experience I came in touch with many other great men to so-called great men but here was a man who was really great not in the conventional sense but in the true sense of the word but and so this year meant terrifically much to me at that time he was already branching out in a new direction called distant parallelism with which I somehow could not could not reconcile myself I didn't think that the solution was in that direction Einstein himself gave up after a few years to give up the whole trend and started something new but on account of the fact that in we couldn't see eye to eye in this question I worked on it on him on a different problem connected with his original gravitational theory which I finished in that year and later published it in 31 so the became became to a friendly departure I mean we it was over it was over but I came in very friendly personal relations with Einstein and I visited him again and again later on when he went to to the States and I was also in the States we correspondent we have a rather lively correspondence and whenever I had a chance to see him in Princeton I made it to a point to visit him there it was basically a lifelong friendship so I came in in again went back to Frankfurt twenty-nine and I was there up to 31 in 1931 the military or the nationalistic agitation was already very strong in Germany and very unpleasant I again personally had no bad experiences but it was a very unpleasant feeling what's going to happen

\subsection{Assistant professor in Purdue University}
at that time again I had a if I would be inclined to superstition I would say head and an archangel who was always always protecting me so when the Hungarian anti-semitism started I came to Germany when the German anti-semitism started I came to the United States ahead of them ahead of the issue because it was in 31 one and a half years before Hitler came to power when I got a call to come to Purdue University where the head of the department who came originally from Vienna and who read the German publications was familiar with my publications and he invited me to come as a visiting professor for one year I accepted that position and I was very happy because I felt that this is the chance of a lifetime I came over to the States with practically no English I was never interested in learning English I didn't take English as a modern language in school I had practically I just after getting after getting that invitation I started to learn frantically some English phrases but it was terribly little little what I knew however I had the great luck that I made good contacts with the faculty and I had one particular friend or on the faculty who was sitting in my lectures because I had to give lectures from the very beginning to a kind of faculty audience who watched all the time made notes and at the end of every lecture pointed out the mistakes I made in the pronunciation and the wrong wrong way of constructing sentences so I learned tremendously in that year I might say that this first year in the States was practically only used up with learning the English because to me it was a question of of existence of survival that I had to learn the English and it's amazing if I look back to this first year that at the same time I published a paper which to my feeling was of considerably interest in the importance in connection with the Einsteinian theory so I had time still to do research but I I slept very little because I was always in a very excited state and in a approximately one year I mean I mean I had the experience that even without any particularly great gift for languages but if you are in a country where you are forced to use the language of the country all the time and I'm I definitely avoided to speak either German or Hungarian with anybody I made it to a point all the time to use English and from the very beginning trying to think in English which after a while I succeeded and then gradually the vocabulary extends and you come to the point where you can first of all understand what the papers say then to the point where you can read books then to the point where you can read newspapers because to read an American newspaper is much more than just reading it an English book they have their own slang and so when you are already able to read the newspaper then you are already a rather advanced and finally to telephone I mean if you can telephone you in a language then you are at home in that language because then you notice but it remembers the distorting device a telephone you have to be quite at home in a language before you are able to do telephone in that language before you are able to understand a fluent conversation in that language so at the end of one year I was at home in English practically at home I knew a great deal I could express myself fluently I could give lectures fluently I would like to say that my success as a lecturer depended very much on the fact that I always prepared my lectures exceedingly carefully although I was able to talk freely after the first two months that I didn't have to read the manuscript but I spent five to six hours for every lecture that I gave at the University on a graduate level they intentionally didn't give me any undergraduate students because they of course would complain immediately for the foreign accent and so on they were treating treating me exceedingly well I made very good contact with the head of the department that mathematics department who was not a creative scientist in any way by the charming man and so this year in a middle middle western university meant a great deal the middle West is a conservative part of the States very conservative outlook in a certain sense we might say narrow-minded but in another respect tremendously good hearted tremendously warm-hearted and tremendously tolerant and at that time there was a great deal of interest in somebody who came from a foreign country and could tell people things how they were in Europe in a completely different world they were very much interested in that I got invitations almost from the very beginning and I was taken full size into their academic life so this was a wonderful experience after one year I made already good contacts and I got it's true it was as a depression it was a depression in in the United States it couldn't give me a full position they gave me half a yet position because I thought I wouldn't like to give up my connections with Germany perhaps I could do it half a year in the States up here in Germany now this was a pipe dream because in the mean time Hitler came to power and he didn't work out but my father was fairly well-to-do he could stand it that I would stay at home for half a year with my parents and then this is what I have done from 1933 on when Hitler came to power and then up to 38 38 it was after Munich no I think I still came home in 39 just before the war breakout just before the world war broke out for the last time in Europe and I actually succeeded in leaving with the last with the last steamer that left share bore to the United States it was perhaps a question of life or death at that time I was able to to save also the life of my son he was at that time a young boy of seven whom I took a long and after that I I got a full appointment and after 38 or 39 I had a full appointment in at Purdue University in the mathematics department I was lecturing mostly to physicists because it was this professor of physics who brought me over but I was appointed in the department of mathematics and I gave a course in in the variational principles of mechanics which became one of my favorite topics they were favorite subjects on which I wrote a book which was a great success and it is still used a great deal it is now in the fourth edition came out in 49 for the first time and the last edition came out only a year ago and the fourth edition and the I learned a great deal myself and

\subsection{Lectures in Hamiltonian dynamics and approximation methods}
the remarkable fact about this lecture course was that whereas usually they were remarkably ingenious methods of Hamilton the famous Irish mathematician where only and only treated as a kind of appendix at the end of the semester at once mistake was I had two semesters at my disposal and I realized the tremendous importance of Hamiltonian methods so that I had a full semester to develop an Hamiltonian dynamics and that again was of very great importance for my little scientific development that I was completely at home in in the Hamiltonian methods now that brings me up now to to the time of the new University and the end of the Purdue University which was in 46 actually in 1938 I made a discovery not in the line of relativity which was my basic field but it so happened that I had to give a course see amazing how much a person learns by giving a course a German colleague of mine an old professor of physics made one's a very witty remark he was altogether a very witty man in Freiburg that you see three times that you don't understand anything of the subject the first time when you learned is a student the second time when you give a lecture course in it in the subject in the third time when you write a book on that so you use you notice three times that you don't understand anything on the subject it's a very good remark so I learned a great deal from a lecture course which I had to give to the students of the physics department on approximation methods it was considered important that they should learn about approximations and I got interested in that that was that was the reason that I have drifted into the field of approximation methods and into a field which one today cause numerical analysis to the extent in 1938 I made a very fundamental discovery in that line which was exceedingly useful for computational purposes a kind of telescoping procedure where a very long a very long stretch of terms could be telescope into a small number of so that you could gain tremendously in computational time I was not interested in how much time you gain but I was interested in the analytical aspects how interesting it is that you can certain procedures exist which which can obtain a certain approximation of a certain accuracy in a few steps where before you needed a tremendously large number of steps I was interested in the analytical aspect but it so happened and when the big electronic machines came and when when during the war it became important to work on certain highly complex military problems these metals were exceedingly practical and I got a call to the National Bureau of Standards to work there as a consultant it that was in the year 43 44 I have trusted a leave of absence from Purdue University and then in 46 a previous student of mine who became the head of research the head of the physical research in the Boeing Airplane Company and he knew of my lectures because he himself took those lectures and of my interest and of my knowledge and expertise in the field of computations and approximation Theory he invited me to join the Boeing Airplane Company as a member of the so-called physical research unit I had the principal in my life always to say yes because I thought it was always interesting perhaps it's a kind of nomad nomad character maybe or maybe there's a kind of gypsy blood in me I don't know but I like changes I like to change the scenery I like to learn new things that to be confronted with new situations whenever somebody says wouldn't you like to come to my place and I said always yes so in 1946 I resigned from Purdue University and went to the Boeing Airplane Company in Seattle where I worked for three years in this physical research unit this was exceedingly interesting time I came in touch with highly educated engineers they were excellent and and years but they didn't understand the formal aspects of mathematics and my problem was to translate a problem from its physical aspects into its mathematical aspects in the form that an engineer could understand it it was a a kind of reinterpretation or just like a an interpreter to take a physical situation translate it into engineering terms and then to take the engineering situation apply the principles of mathematics get the solution and going back again to the physical situation in this I was successful and they I liked I liked to this contact with these highly trained and excellent a highly intelligent engineers they were exceedingly narrow minded in certain respects had absolutely no feeling for labor problems they always lined up with the capitalists and absolutely no feeling that there are poor people and that they are also they also have certain rights to live it was quite astonishing how narrow their social outlook was but in their own field they were excellent

\subsection{Institute for Numerical Analysis in Los Angeles and the McCarthy Era}
and then in 49 again I got a call to the newly at that time I published I had publications already in particularly this 38 paper with so-called chebyshev polynomials which I showed are exceedingly useful in numerical analysis I got a call to the newly founded Institute for numerical analysis in Los Angeles again I I accepted it immediately and I had four years three or four years of very interesting and wonderful years in California that was in 49 up to 53 in 53 again an unpleasant upheaval happened on account of the so-called McCarthy era where everybody who had some liberal meanings leanings anybody with liberal leanings had immediately he was immediately called a communist now the Institute for which I worked blew up automatically it was under there it was under the sponsorship of the National Bureau of Standards and for some reasons this natural Bureau of Standards was immediately attacked as vie by the McCarthy eyes as a hotbed of homosexuals and and communists so what these two things have to do with each other it's a mystery to me but that is what they claimed so my chief immediately lost his position and the Institute came into a very precarious position

\subsection{Dublin Institute for Advanced Studies}
at that time I by accident I came in correspondence with the Dublin Institute for Advanced Studies this EE Institute for Advanced Studies was an FIR dinner was there the famous physicist was the director of that Institute at that time and it so happened that somebody who was an expert in editing the papers of Hamilton the scientific papers are of Hamilton died and he thought that perhaps since I worked so much in Hamiltonian dynamics perhaps I would be the right person to take over his his job in editing the Hamiltonian papers which have been published by the Royal Irish Academy so I accepted that I got a leave of absence from the National Bureau standards and left for Ireland now Ireland with its old-world atmosphere appeal to me tremendously it was an entirely new world but it was somehow like a world to which I came back and I liked it very much and I came in good relations with the staff of the Institute I after one year of visiting professorship I did come back to the States for half a year I worked for the North American Aviation corporation in some problems of aerodynamics and numerical analysis but then president of alehrer who at that time was not president but he was a minister the prime prime minister he he offered me and he offered me a position a permanent position a senior position at the Institute for Advanced Studies this to me was very great not only a great honor but it was also of great practical advantage I was already 60 years old in the States I would have had to resign to in in about five years I would have to to go into retirement the 65 is a retirement age moreover I worked for industry for a number of years whereas it's very easy to go from the academic life into industry it's not so easy to go back from industry in the academic life and I wanted to ahead a great longing to go back again to give up these applied applied problems and numerical problems and computational problems and these applications of engineering mathematics to engineering problems and explaining to engineers what they have to do I thought it is enough and I wanted to go back to more fundamental problems again I had a great longing to return to my first love which was a lot EVT and to do again research in that life so this offer came to me like a real salvation because at our Institute in Dublin one can work up to 75 which is unusual so I still had 15 years of active academic life wherever our Institute in Dublin is organ along remarkably liberalize it is a pure Research Institute we don't have to give courses we don't have under graduates we have only post doctoral we have we have only postdoctoral students as we call them scholars who work at the Institute under very free conditions they can line up this one usually line up is one of the senior professors work under him and then they leave after one or maximum two years we don't grant any degrees they are mostly students who have already their PhD degree we make only exception with Irish students because they have so little chances to make their PhD so they can work on a certain research subject which they can later on submit as a PhD but otherwise we are free of any any thing with which one usually is bothered in a in a university industry but here we have it's a purely scientific institution and I think it is remarkable that a small country like Ireland can allows itself such a luxury that was president de Valera's great interest in mathematics he himself was a mathematician who went into politics and so he organized this Institute along exceedingly pleasant lines and I never regretted the fact that I went to Ireland I thought when I finish with Ireland maybe I can go back to the states and I was in the States many times for visits but then my health broke down and I could not realize these plans and now that brings in it brings me to the to the present day I'm fortunately still alive and I continue my my work as before I have my home in the Institute and I don't see any particular change from professor senior professor to emeritus professor I don't see any particular trainer I'm on very good terms with the director and with everybody else and enjoy the fact that I can still work and now I am here in Manchester

\subsection{Consultant at National Bureau Standards and researcher at Boeing}
and I got a call I got an invitation from professor Hopkins and dr. Butler which I enjoyed tremendously to give six lectures on Einstein I gave already six lectures on the basis of a book that I wrote in no actually that book was based on these six lectures which I gave in Ann Arbor University I don't repeat this like I in fact I didn't look even into the book because I want to show some other aspects of them of the problems but I have a chance to talk on Einstein for on six occasions and now this time I had a chance to talk about myself but I think you noticed that I didn't want just to point out how many wonderful things I have done in life but perhaps it was more interesting to to give to give the survey over of a scientific career which brought me into many countries in in into many situations interesting situations I had a rich and interesting life which I always appreciated tremendously and now I am here and you had you had the kindness to take this videotape of my life so I thank you very much for your interest and for your patience


\section{About Mathematics}
\subsection{Introduction}
\subsection{Pure and Applied Mathematics}
\subsection{Approximation theory, Otto Szász, and Chebyshev polynomials}
\subsection{Economization of power series and the Lanczos $\tau$ method}
\subsection{National Bureau of Standards}
\subsection{Singular matrices and the SVD}
\subsection{Boeing and the Lanczos algorithm}
\subsection{Importance of Fourier analysis}
\subsection{Gauss and hypergeometric series}
\subsection{How to lecture and write books}
\subsection{Lanczos' favourite mathematicians}
\subsection{The right mathematical notation}
\subsection{History of mathematics, education, and deep understanding}

\section{About Albert Einstein}
\subsection{Introduction}
\subsection{The young Albert Einstein in 190}
\subsection{Position at the Swiss Patent Office}
\subsection{The 1905 paper on special relativity}
\subsection{Work of Poincaré and Lorenz}
\subsection{Einstein's style of thinking}
\subsection{Reception of the papers in 1909, positions in Zurich, Prague, Berlin}
\subsection{The 1905 paper on Brownian motion}
\subsection{Work on the photoelectric effect}
\subsection{$E=mc^2$}
\subsection{God doesn't play dice}
\subsection{Einstein's mathematical limitations}
\subsection{Gravitation and general relativity}
\subsection{Einstein as a pacifist}

\end{document}
